% !TEX root = ../main-es.tex
% Las cajas de ayuda no son contenido: texcount las ignora (ver `make wordcount`).
%TC:macro \limite [ignore,ignore]
%TC:macro \nota [ignore]
%TC:macro \TODO [ignore]
%TC:macro \parencite [ignore]
%TC:macro \textcite [ignore]
\section{Contexto del proyecto}
\limite{0}{Contexto del proyecto: describa brevemente la motivación del proyecto, el
problema que busca abordar y su relevancia.}

\subsection*{Motivación}

El proyecto nace de un sistema real: \textbf{OASys}, un motor de ejecución agéntica escrito en el que un agente basado en inteligencia artificial generativa (LLM/SOM) puede \emph{leer y reescribir el programa que lo está ejecutando}, mediante las herramientas vinculadas al \texttt{filesystem}. La motivación de diseño es sostener tareas de horizonte largo: Procesos que ni un LLM suelto completa (causado por degradación de contexto) ni un flujo determinista afronta (el entorno cambia y el flujo es estático). Un agente que reescribe su propio flujo escapa, en principio, a ambas limitaciones.


El problema es que esa capacidad nunca se midió. Hoy el sistema resuelve sus tres decisiones centrales por intuición: el modo de aplicación inmediata de ediciones (\texttt{run=hot}) está declarado pero no implementado, el número máximo de iteraciones de deliberación es una constante fijada a mano, y el único freno contra reescrituras improductivas es la aprobación de un humano.

\subsection*{Problema}

Al leer la taxonomía clásica de agentes con OASys en la cabeza \parencite{russellNorvig2020}, el sistema completo resulta estar descrito en la literatura con otro vocabulario: el programa que ejecuta el agente es un plan jerárquico, leerlo es monitoreo de ejecución, reescribirlo es replanificación, y aplicar las ediciones en frío o en caliente es planeación de lazo abierto contra lazo cerrado.

Si el problema es clásico, cabe la pregunta clásica: replanificar cuesta (deliberar consume tiempo durante el cual el mundo sigue cambiando) y su beneficio crece con el dinamismo del entorno. Beneficio creciente contra costo aproximadamente constante implica un punto de cruce. La pregunta del proyecto es:

\begin{quote}
  \proyPregunta
\end{quote}

El problema se modela como un sistema multiagente en MESA \parencite{kazil2020mesa}: Un
agente recolector con seis arquitecturas (de reflejo simple a lazo cerrado con
gobernanza) que opera en una retícula con recursos móviles, donde el replanificador es
un oráculo abstracto con competencia $p$ y costo $c$ por edición.

\subsection*{Relevancia}

La relevancia implica hacia el sistema real que la frontera medida entre lazo abierto y replanificación es el argumento cuantitativo que falta para decidir si vale la pena su uso, con qué presupuesto de ediciones y con qué modelo de costo. Hacia la teoría agéntica, el proyecto usa su contenido como pieza funcional: El planificador que el agente invoca en cada reescritura es A*, y la comparación entre arquitecturas es exactamente la taxonomía de agentes y entornos del curso.
